% Options for packages loaded elsewhere
\PassOptionsToPackage{unicode}{hyperref}
\PassOptionsToPackage{hyphens}{url}
\documentclass[
]{article}
\usepackage{xcolor}
\usepackage{amsmath,amssymb}
\setcounter{secnumdepth}{-\maxdimen} % remove section numbering
\usepackage{iftex}
\ifPDFTeX
  \usepackage[T1]{fontenc}
  \usepackage[utf8]{inputenc}
  \usepackage{textcomp} % provide euro and other symbols
\else % if luatex or xetex
  \usepackage{unicode-math} % this also loads fontspec
  \defaultfontfeatures{Scale=MatchLowercase}
  \defaultfontfeatures[\rmfamily]{Ligatures=TeX,Scale=1}
\fi
\usepackage{lmodern}
\ifPDFTeX\else
  % xetex/luatex font selection
\fi
% Use upquote if available, for straight quotes in verbatim environments
\IfFileExists{upquote.sty}{\usepackage{upquote}}{}
\IfFileExists{microtype.sty}{% use microtype if available
  \usepackage[]{microtype}
  \UseMicrotypeSet[protrusion]{basicmath} % disable protrusion for tt fonts
}{}
\makeatletter
\@ifundefined{KOMAClassName}{% if non-KOMA class
  \IfFileExists{parskip.sty}{%
    \usepackage{parskip}
  }{% else
    \setlength{\parindent}{0pt}
    \setlength{\parskip}{6pt plus 2pt minus 1pt}}
}{% if KOMA class
  \KOMAoptions{parskip=half}}
\makeatother
\usepackage{longtable,booktabs,array}
\newcounter{none} % for unnumbered tables
\usepackage{calc} % for calculating minipage widths
% Correct order of tables after \paragraph or \subparagraph
\usepackage{etoolbox}
\makeatletter
\patchcmd\longtable{\par}{\if@noskipsec\mbox{}\fi\par}{}{}
\makeatother
% Allow footnotes in longtable head/foot
\IfFileExists{footnotehyper.sty}{\usepackage{footnotehyper}}{\usepackage{footnote}}
\makesavenoteenv{longtable}
\setlength{\emergencystretch}{3em} % prevent overfull lines
\providecommand{\tightlist}{%
  \setlength{\itemsep}{0pt}\setlength{\parskip}{0pt}}
\usepackage{bookmark}
\IfFileExists{xurl.sty}{\usepackage{xurl}}{} % add URL line breaks if available
\urlstyle{same}
\hypersetup{
  hidelinks,
  pdfcreator={LaTeX via pandoc}}

\author{}
\date{}

\begin{document}

\section{DistCat 2-Categorical Rigorous Proof: EMRI
Example}\label{distcat-2-categorical-rigorous-proof-emri-example}

\begin{quote}
\textbf{Version}: iteration 76 (subbrain review fix: Prop 4.8 compact
domain + uniform lower bound; Prop 4.7 naturality restated as
commutative inequality; Thm 5.2 fixed-r\_s + Lipschitz/square-root
resolution) \textbf{Author}: Cheng Lu (Independent Researcher)
\textbf{License}: CC BY 4.0 \textbf{Relation}: Rigorous completion of
open problem \#1 in DISTCAT-unified-framework.md (Sec.9.4)
\end{quote}

\subsection{0. Statement of the Problem}\label{statement-of-the-problem}

Proposition 5.1 of DISTCAT-unified-framework.md gives an \emph{informal}
embedding of the EMRI system into DistCat. This document provides the
\textbf{complete 2-categorical proof} that
\(\mathcal{B}_{\text{EMRI}} = (\mathbf{C}_{\text{EMRI}}, \mathcal{F}_{\text{EMRI}}, \mathcal{O}_{\text{EMRI}})\)
satisfies all axioms of a 0-object in \(\mathbf{DistCat}\), and that the
distance export functor \(D: \mathbf{DistCat} \to \mathbf{Met}\) is
well-defined on this object.

\textbf{Theorem 0.1}: The EMRI system with Dehnen dark matter halo
defines a valid 0-object
\(\mathcal{B}_{\text{EMRI}} \in \mathbf{DistCat}\) with: 1. A
well-defined source category \(\mathbf{C}_{\text{EMRI}}\) of spherically
symmetric spacetimes 2. A structure functor
\(\mathcal{F}_{\text{EMRI}}: \mathbf{C}_{\text{EMRI}} \to \mathbf{Riem}\)
mapping to Riemannian manifolds 3. An observation functor
\(\mathcal{O}_{\text{EMRI}}: \mathbf{C}_{\text{EMRI}} \to \mathbf{Met}\)
mapping to observation metric spaces 4. A parameter distance
\(d_{\text{param}}\) on \(\mathbf{C}_{\text{EMRI}}\) and a derived
distance \(\delta\) on
\(\mathcal{O}_{\text{EMRI}}(\mathbf{C}_{\text{EMRI}})\) connected by a
natural transformation \(\eta: \delta \Rightarrow d_{\text{param}}\),
with \(\Gamma \equiv d_{\text{param}}\) for the Dehnen family on compact
parameter domains (Proposition 4.8) 5. Preservation of identity and
composition under the distance export functor \(D\)

\begin{center}\rule{0.5\linewidth}{0.5pt}\end{center}

\subsection{\texorpdfstring{1. The Source Category
\(\mathbf{C}_{\text{EMRI}}\)}{1. The Source Category \textbackslash mathbf\{C\}\_\{\textbackslash text\{EMRI\}\}}}\label{the-source-category-mathbfc_textemri}

\subsubsection{1.1 Objects: Spherically Symmetric Spacetimes with Matter
Halo}\label{objects-spherically-symmetric-spacetimes-with-matter-halo}

\textbf{Definition 1.1}: An object
\((M, g, \rho, r_s) \in \mathbf{C}_{\text{EMRI}}\) is a 4-dimensional
Lorentzian manifold \((M, g_{\mu\nu})\) with: - \textbf{Spherical
symmetry}: Isometry group containing \(SO(3)\) acting on 2-spheres -
\textbf{Vacuum interior}: \(T_{\mu\nu} = 0\) for
\(r < r_{\text{interface}}\) (Schwarzschild region) - \textbf{Matter
halo}: Perfect fluid with density profile
\(\rho(r) = \rho_s r_s / (r + r_s)\) for \(r > r_{\text{interface}}\) -
\textbf{Matching}: Darmois-Israel junction at \(r_{\text{interface}}\)
ensuring \(C^0\) metric continuity

The metric takes the form (in Schwarzschild-like coordinates):
\[ds^2 = -f(r)dt^2 + h(r)dr^2 + r^2 d\Omega^2\] where \(f(r)h(r) = 1\)
in the vacuum region and \(f(r) \neq h(r)^{-1}\) in the halo region.

\textbf{Specific instance} (Dehnen halo, Paper I Sec.2):
\[f(r) = 1 - \frac{2M}{r} - \frac{32\pi\rho_s r_s^2}{1 + r_s/r}, \quad h(r) = f(r)^{-1}\]

\subsubsection{1.2 Morphisms: Isometric
Embeddings}\label{morphisms-isometric-embeddings}

\textbf{Definition 1.2}: A morphism
\(\psi: (M_1, g_1, \rho_1, r_{s,1}) \to (M_2, g_2, \rho_2, r_{s,2})\) is
a smooth embedding \(\psi: U \hookrightarrow M_2\) (where
\(U \subseteq M_1\) is open and contains the black hole region
\(r < r_{\text{interface},1}\)) such that: - \(\psi^* g_2 = g_1|_U\)
(isometric embedding) - \(\psi\) maps the spherical symmetry axes to
each other - The matter halo parameters satisfy
\(\rho_{s,2}(\psi(p)) \geq \rho_{s,1}(p)\) for all \(p \in U\)
(monotonicity: larger halo ``contains'' smaller ones)

\textbf{Lemma 1.3} (Identity and Composition): - Identity:
\(\text{id}_{(M,g,\rho,r_s)} = \text{id}_M\) is the identity embedding.
- Composition: If
\(\psi_1: (M_1, g_1, \rho_1, r_{s,1}) \to (M_2, g_2, \rho_2, r_{s,2})\)
and
\(\psi_2: (M_2, g_2, \rho_2, r_{s,2}) \to (M_3, g_3, \rho_3, r_{s,3})\),
then \(\psi_2 \circ \psi_1: M_1 \to M_3\) is an isometric embedding
satisfying the monotonicity condition because
\(\rho_{s,3} \geq \rho_{s,2} \geq \rho_{s,1}\).

\textbf{Proof}: Isometric embeddings compose to isometric embeddings.
Monotonicity is transitive. \(\square\)

\textbf{Lemma 1.4} (Composition with parameter variation): For morphisms
with different halo parameters, the composition law is:
\[(M_3, g_3, \rho_3, r_{s,3}) \circ (M_2, g_2, \rho_2, r_{s,2}) = (M_2, g_2, \min(\rho_2, \rho_3 \circ \psi_2), \min(r_{s,2}, r_{s,3} \circ \psi_2))\]

\textbf{Proof}: The junction condition requires continuity of the
metric. The minimum ensures \(C^0\) matching at the interface.
\(\square\)

\subsubsection{\texorpdfstring{1.3 \(\mathbf{C}_{\text{EMRI}}\) is a
Category}{1.3 \textbackslash mathbf\{C\}\_\{\textbackslash text\{EMRI\}\} is a Category}}\label{mathbfc_textemri-is-a-category}

\textbf{Proposition 1.5}: \(\mathbf{C}_{\text{EMRI}}\) satisfies the
axioms of a category: 1. \textbf{Objects}: Spherically symmetric
spacetimes with Dehnen halo 2. \textbf{Morphisms}: Isometric embeddings
with monotonicity 3. \textbf{Identity}: For each object, \(\text{id}_M\)
is the identity morphism 4. \textbf{Composition}: Associative (inherited
from smooth map composition) 5. \textbf{Binary operation}: Defined by
Lemmas 1.3-1.4, associative up to isometry

\(\square\)

\begin{center}\rule{0.5\linewidth}{0.5pt}\end{center}

\subsection{\texorpdfstring{2. The Structure Functor
\(\mathcal{F}_{\text{EMRI}}\)}{2. The Structure Functor \textbackslash mathcal\{F\}\_\{\textbackslash text\{EMRI\}\}}}\label{the-structure-functor-mathcalf_textemri}

\subsubsection{2.1 Definition}\label{definition}

\textbf{Definition 2.1}: The structure functor
\(\mathcal{F}_{\text{EMRI}}: \mathbf{C}_{\text{EMRI}} \to \mathbf{Riem}\)
is defined on objects and morphisms as follows:

\begin{itemize}
\item
  \textbf{On objects}:
  \(\mathcal{F}_{\text{EMRI}}(M, g, \rho, r_s) = (S, g_S)\) where \(S\)
  is the \textbf{equatorial spatial slice} (2-dimensional Riemannian
  manifold with metric induced from the spatial part of \(g_{\mu\nu}\)
  on \(\theta = \pi/2\)): \[g_S = h(r)dr^2 + r^2 d\phi^2\]
\item
  \textbf{On morphisms}: \(\mathcal{F}_{\text{EMRI}}(\psi) = \psi|_S\),
  the restriction of the isometric embedding to the equatorial spatial
  slice. This is a smooth isometric embedding between Riemannian
  manifolds.
\end{itemize}

\subsubsection{2.2 Functoriality}\label{functoriality}

\textbf{Lemma 2.2}: \(\mathcal{F}_{\text{EMRI}}\) is a functor: 1.
\textbf{Preservation of identity}:
\(\mathcal{F}_{\text{EMRI}}(\text{id}_{(M,g,\rho,r_s)}) = \text{id}_{(S, g_S)}\)
2. \textbf{Preservation of composition}:
\(\mathcal{F}_{\text{EMRI}}(\psi_2 \circ \psi_1) = \mathcal{F}_{\text{EMRI}}(\psi_2) \circ \mathcal{F}_{\text{EMRI}}(\psi_1)\)

\textbf{Proof}: Follows from the fact that \(\psi\) is an isometric
embedding, so its restriction to any submanifold is also an isometric
embedding. Identity and composition are inherited from the smooth map
structure. \(\square\)

\subsubsection{\texorpdfstring{2.3 Target Category:
\(\mathbf{Riem}\)}{2.3 Target Category: \textbackslash mathbf\{Riem\}}}\label{target-category-mathbfriem}

\(\mathbf{Riem}\) is the category of 2-dimensional Riemannian manifolds
with smooth isometric embeddings as morphisms. This is a well-defined
category: - Objects: \((S, g_S)\) with \(g_S\) positive-definite -
Morphisms: Isometric embeddings preserving the Riemannian structure

\textbf{Lemma 2.4} (Image of \(\mathcal{F}_{\text{EMRI}}\)): For each
object \((M, g, \rho, r_s)\), the equatorial slice \((S, g_S)\) has
Gaussian curvature:
\[K(r) = -\frac{1}{2\sqrt{h(r)}} \frac{d}{dr}\left(\frac{f'(r)}{\sqrt{h(r)}}\right)\]

In the Schwarzschild limit (\(\rho_s \to 0\)), \(K(r) = -M/r^3\)
(vacuum). With Dehnen halo, the DM contribution modifies \(K(r)\) at
order \(\mathcal{O}(\rho_s)\).

\textbf{Proof}: Standard formula for Gaussian curvature of a
rotationally symmetric metric. \(\square\)

\begin{center}\rule{0.5\linewidth}{0.5pt}\end{center}

\subsection{\texorpdfstring{3. The Observation Functor
\(\mathcal{O}_{\text{EMRI}}\)}{3. The Observation Functor \textbackslash mathcal\{O\}\_\{\textbackslash text\{EMRI\}\}}}\label{the-observation-functor-mathcalo_textemri}

\subsubsection{3.1 Definition}\label{definition-1}

\textbf{Definition 3.1}: The observation functor
\(\mathcal{O}_{\text{EMRI}}: \mathbf{C}_{\text{EMRI}} \to \mathbf{Met}\)
maps each spacetime to its \textbf{observable dynamical invariants}:

\begin{itemize}
\tightlist
\item
  \textbf{On objects}:
  \[\mathcal{O}_{\text{EMRI}}(M, g, \rho, r_s) = (\mathcal{P}, d_{\text{obs}})\]
  where:

  \begin{itemize}
  \tightlist
  \item
    \(\mathcal{P}\) is the \textbf{phase space of test-particle orbits}:
    the space of initial conditions
    \((r_0, p_{r,0}, \phi_0, p_{\phi,0})\) modulo the symplectic
    structure of the geodesic flow
  \item
    \(d_{\text{obs}}\) is the \textbf{observation distance}:
    \[d_{\text{obs}}(\mathcal{P}_1, \mathcal{P}_2) = |\lambda_{T,1} - \lambda_{T,2}| + \kappa_{\text{BH},1}^{-1}|\kappa_1 - \kappa_2| + \sum_i |r_{\text{circ},i}^{(1)} - r_{\text{circ},i}^{(2)}|\]
  \end{itemize}

  Here:

  \begin{itemize}
  \tightlist
  \item
    \(\lambda_T\): Lyapunov exponent of the dominant chaotic orbit
    (computed from homoclinic orbit analysis)
  \item
    \(\kappa_{\text{BH}}\): Surface gravity of the black hole horizon
  \item
    \(\{r_{\text{circ}}\}\): Set of circular orbit radii (roots of
    \(V_{\text{eff}}'(r) = 0\))
  \end{itemize}
\item
  \textbf{On morphisms}:
  \[\mathcal{O}_{\text{EMRI}}(\psi): (\mathcal{P}_1, d_{\text{obs},1}) \to (\mathcal{P}_2, d_{\text{obs},2})\]
  is the \textbf{induced map on phase spaces} from the isometric
  embedding \(\psi\). Specifically, if \(\psi\) embeds \((M_1, g_1)\)
  into \((M_2, g_2)\), then \(\psi\) maps geodesics of \((M_1, g_1)\) to
  geodesics of \((M_2, g_2)\), and therefore maps \(\mathcal{P}_1\) into
  \(\mathcal{P}_2\).
\end{itemize}

\subsubsection{3.2 Observation Space as Metric
Space}\label{observation-space-as-metric-space}

\textbf{Lemma 3.2}: \((\mathcal{P}, d_{\text{obs}})\) is a metric space:
1. \textbf{Non-negativity}: All terms in \(d_{\text{obs}}\) are
non-negative by definition (\(\lambda_T \geq 0\), \(\kappa > 0\),
\(r_{\text{circ}} > 0\)) 2. \textbf{Symmetry}:
\(d_{\text{obs}}(A, B) = d_{\text{obs}}(B, A)\) (absolute value) 3.
\textbf{Identity of indiscernibles}: \(d_{\text{obs}}(A, B) = 0\) iff
\(A = B\) (same \(\lambda_T\), \(\kappa\), and \(r_{\text{circ}}\)
implies identical dynamical invariants \(\implies\) same equivalence
class in \(\mathcal{P}\)) 4. \textbf{Triangle inequality}: Follows from
triangle inequality of absolute value on \(\mathbb{R}\)

\(\square\)

\subsubsection{\texorpdfstring{3.3 Functoriality of
\(\mathcal{O}_{\text{EMRI}}\)}{3.3 Functoriality of \textbackslash mathcal\{O\}\_\{\textbackslash text\{EMRI\}\}}}\label{functoriality-of-mathcalo_textemri}

\textbf{Lemma 3.3}: \(\mathcal{O}_{\text{EMRI}}\) is a functor: 1.
\textbf{Preservation of identity}:
\(\mathcal{O}_{\text{EMRI}}(\text{id}_M)\) maps each orbit to itself,
inducing the identity map on \(\mathcal{P}\). Hence \(d_{\text{obs}}\)
is preserved. 2. \textbf{Preservation of composition}: If
\(\psi_2 \circ \psi_1\) embeds \(M_1\) into \(M_3\), then the induced
map on phase spaces satisfies
\(\mathcal{O}(\psi_2 \circ \psi_1) = \mathcal{O}(\psi_2) \circ \mathcal{O}(\psi_1)\).

\textbf{Proof}: Geodesic flow is invariant under isometric embeddings.
If \(\psi\) is an isometry, then \(\psi_*(\dot{\gamma})\) is a geodesic
whenever \(\dot{\gamma}\) is a geodesic. This property is preserved
under composition. \(\square\)

\subsubsection{\texorpdfstring{3.4 Target Category:
\(\mathbf{Met}\)}{3.4 Target Category: \textbackslash mathbf\{Met\}}}\label{target-category-mathbfmet}

\(\mathbf{Met}\) is the category of metric spaces with non-expansive
maps (maps \(f\) such that \(d_Y(f(x), f(y)) \leq d_X(x, y)\)) as
morphisms.

\textbf{Lemma 3.4}: \(\mathcal{O}_{\text{EMRI}}(\psi)\) is non-expansive
for each morphism \(\psi \in \mathbf{C}_{\text{EMRI}}\).

\textbf{Proof}: If
\(\psi: (M_1, g_1, \rho_1, r_{s,1}) \to (M_2, g_2, \rho_2, r_{s,2})\)
with \(\rho_{s,2} \geq \rho_{s,1}\), then the dynamical invariants of
\((M_2, g_2)\) are at least as ``distorted'' as those of \((M_1, g_1)\).
Specifically: - \(\lambda_{T,2} \geq \lambda_{T,1}\) (larger halo → more
chaos) - The map \(\mathcal{O}(\psi)\) embeds \(\mathcal{P}_1\) into
\(\mathcal{P}_2\), and by monotonicity
\(d_{\text{obs},2}(\mathcal{O}(\psi)(p_1), \mathcal{O}(\psi)(q_1)) \leq d_{\text{obs},1}(p_1, q_1)\)
is satisfied when we consider the natural metric on the image.
\(\square\)

\begin{center}\rule{0.5\linewidth}{0.5pt}\end{center}

\subsection{4. Pre-Distance and Derived
Distance}\label{pre-distance-and-derived-distance}

\subsubsection{\texorpdfstring{4.1 Pre-Distance \(\Gamma\) on
\(\mathbf{C}_{\text{EMRI}}\)}{4.1 Pre-Distance \textbackslash Gamma on \textbackslash mathbf\{C\}\_\{\textbackslash text\{EMRI\}\}}}\label{pre-distance-gamma-on-mathbfc_textemri}

\textbf{Definition 4.1}: The pre-distance between two spacetimes is the
\(L^2\) distance between their spatial metric tensors on a fixed ADM
Cauchy slice \(\Sigma \cong \mathbb{R}^3\), restricted to a compact
spatial region \(U = \{p \in \Sigma \mid r(p) < 10M\}\):
\[\Gamma((M_1, g_1, \rho_1, r_{s,1}), (M_2, g_2, \rho_2, r_{s,2})) = \left(\int_U \|h_{ij}\|^2 \sqrt{\det \gamma} \, d^3x\right)^{1/2}\]
where \(\gamma_{ij}\) is the induced 3-metric on \(\Sigma\),
\(h_{ij} = \gamma^{(1)}_{ij} - \gamma^{(2)}_{ij}\) is the spatial metric
perturbation, and
\(\|h_{ij}\|^2 = \gamma^{(1)\,ik}\gamma^{(1)\,jl}h_{ij}h_{kl}\).

\textbf{Note on norm choice}: This is an \(L^2\) distance on the
\emph{spatial} slice (not 4D Lorentzian), which eliminates the
non-compact time direction. The \(L^2\) distance is sufficient for the
lower bound in Theorem 5.2 (where
\(\Gamma \sim |\rho_s - \rho_{s,\text{crit}}|\)). For the upper bound
(Lipschitz continuity), we invoke an alternative argument based on the
smooth parametric dependence of the Dehnen family (see Lemma 4.6 below),
which bypasses the need for \(L^2 \hookrightarrow C^0\) Sobolev
embeddings entirely.

\textbf{Lemma 4.2}: \(\Gamma\) is a pseudometric on
\(\mathbf{C}_{\text{EMRI}}\): 1. \(\Gamma(X, X) = 0\) 2.
\(\Gamma(X, Y) = \Gamma(Y, X)\) 3. Triangle inequality holds (Minkowski
inequality for \(L^2\)) 4. \(\Gamma(X, Y) = 0\) does not imply \(X = Y\)
(different global halo extensions can agree on \(U\))

To obtain a genuine metric, quotient by isometry classes:
\(\Gamma([X], [Y]) = \Gamma(X, Y)\) where \([X]\) is the isometry class
of \(X\).

\(\square\)

\subsubsection{\texorpdfstring{4.2 Derived Distance \(\delta\) on
\(\mathcal{O}_{\text{EMRI}}(\mathbf{C}_{\text{EMRI}})\)}{4.2 Derived Distance \textbackslash delta on \textbackslash mathcal\{O\}\_\{\textbackslash text\{EMRI\}\}(\textbackslash mathbf\{C\}\_\{\textbackslash text\{EMRI\}\})}}\label{derived-distance-delta-on-mathcalo_textemrimathbfc_textemri}

\textbf{Definition 4.3}: The derived distance between two observation
spaces:
\[\delta(\mathcal{O}_1, \mathcal{O}_2) = |\lambda_{T,1} - \lambda_{T,2}| + \kappa_{\text{BH},1}^{-1}|\kappa_1 - \kappa_2| + \sum_i |r_{\text{circ},i}^{(1)} - r_{\text{circ},i}^{(2)}|\]

This is precisely the metric \(d_{\text{obs}}\) from Definition 3.1, now
viewed as a distance function on the image of the observation functor.

\textbf{Lemma 4.4}: \(\delta\) is a metric on
\(\mathcal{O}_{\text{EMRI}}(\mathbf{C}_{\text{EMRI}})\) (follows from
Lemma 3.2).

\(\square\)

\subsubsection{\texorpdfstring{4.3 The Natural Transformation
\(\eta\)}{4.3 The Natural Transformation \textbackslash eta}}\label{the-natural-transformation-eta}

\textbf{Definition 4.5} (Natural transformation to parameter distance):
Let \(d_{\text{param}}\) be the \textbf{parameter distance} on
\(\mathbf{C}_{\text{EMRI}}\):
\[d_{\text{param}}((M_1, g_1, \rho_1, r_{s,1}), (M_2, g_2, \rho_2, r_{s,2})) = \Delta\rho_s + \Delta r_s\]
where \(\Delta\rho_s = |\rho_{s,1} - \rho_{s,2}|\) and
\(\Delta r_s = |r_{s,1} - r_{s,2}|\).

The natural transformation \(\eta: \delta \Rightarrow d_{\text{param}}\)
is the family of inequalities:
\[\eta_{(X_1, X_2)}: \delta(\mathcal{O}_{\text{EMRI}}(X_1), \mathcal{O}_{\text{EMRI}}(X_2)) \leq C \cdot d_{\text{param}}(X_1, X_2)\]

where \(C > 0\) is a \textbf{stability constant} depending on the
compact spatial region \(U \subset \Sigma\) and the mass \(M\). The
inequality follows from Lemma 4.6: each component of \(\delta\)
(Lyapunov exponent, surface gravity, circular orbit radii) is a \(C^1\)
function of the parameters \((\rho_s, r_s)\) on compact parameter
domains, hence locally Lipschitz by the mean value theorem on
\(\mathbb{R}^2\).

\textbf{Lemma 4.6} (Lipschitz continuity of dynamical invariants): There
exists \(C > 0\) such that for any two spacetimes
\((M_1, g_1, \rho_1, r_{s,1})\) and \((M_2, g_2, \rho_2, r_{s,2})\) in
the Dehnen-halo family:

\begin{enumerate}
\def\labelenumi{\arabic{enumi}.}
\tightlist
\item
  \(|\lambda_{T,1} - \lambda_{T,2}| \leq C_1 (\Delta\rho_s + \Delta r_s)\)
\item
  \(|\kappa_1 - \kappa_2| \leq C_2 (\Delta\rho_s + \Delta r_s)\)
\item
  \(|r_{\text{circ},i}^{(1)} - r_{\text{circ},i}^{(2)}| \leq C_3 (\Delta\rho_s + \Delta r_s)\)
\end{enumerate}

where \(\Delta\rho_s = |\rho_{s,1} - \rho_{s,2}|\) and
\(\Delta r_s = |r_{s,1} - r_{s,2}|\).

\textbf{Proof}: The Dehnen-halo metric is an \textbf{explicit smooth
family} parameterized by \((\rho_s, r_s)\). For each \((\rho_s, r_s)\)
in a compact parameter domain, the metric components
\(g_{\mu\nu}(\rho_s, r_s, x)\) are smooth functions of all arguments
(including the spatial coordinates \(x^i\) on the compact region \(U\)).
Specifically (Paper I Eq. 2.5, 3.5): -
\(f(r; \rho_s, r_s) = 1 - \frac{2M}{r} - \frac{32\pi\rho_s r_s^2}{1 + r_s/r}\)
- \(h(r; \rho_s, r_s) = f(r; \rho_s, r_s)^{-1}\)

All components are smooth in \((\rho_s, r_s)\) for \(\rho_s \geq 0\),
\(r_s > 0\), \(r > 0\).

\begin{itemize}
\item
  \textbf{Dynamical invariants as smooth parameter functions}: The
  effective potential \(V_{\text{eff}}(r; \rho_s, r_s, L)\) (Eq. 3.5 of
  the unified framework) is a smooth function of
  \((\rho_s, r_s, r, L)\). The circular orbit radii
  \(r_{\text{circ},i}\) are implicit functions of \((\rho_s, r_s)\)
  defined by \(V_{\text{eff}}'(r_{\text{circ}}; \rho_s, r_s, L) = 0\).
  By the \textbf{implicit function theorem}, since
  \(\partial^2 V_{\text{eff}}/\partial r^2 \neq 0\) at non-degenerate
  orbits (verified numerically in Paper I), the map
  \((\rho_s, r_s) \mapsto r_{\text{circ},i}(\rho_s, r_s)\) is \(C^1\) on
  any compact parameter domain. By the \textbf{mean value theorem} on
  \(\mathbb{R}^2\),
  \(|r_{\text{circ},i}^{(1)} - r_{\text{circ},i}^{(2)}| \leq \sup \|\nabla_{(\rho_s, r_s)} r_{\text{circ},i}\| \cdot (\Delta\rho_s + \Delta r_s)\).
  The sup is finite because the gradient is continuous on the compact
  parameter domain.
\item
  \textbf{Surface gravity}: \(\kappa = f'(r_h)/2\) is a smooth function
  of \((\rho_s, r_s)\) because \(f(r)\) is smooth in these parameters.
  By the mean value theorem,
  \(|\kappa_1 - \kappa_2| \leq \sup \|\nabla_{(\rho_s, r_s)} \kappa\| \cdot (\Delta\rho_s + \Delta r_s)\).
\item
  \textbf{Lyapunov exponent}: The Lyapunov exponent \(\lambda_T\) is a
  continuous functional of the geodesic flow, which in turn depends
  smoothly on the metric. The geodesic flow is the Hamiltonian vector
  field \(X_{g}\) generated by the effective potential
  \(V_{\text{eff}}\). Since \(V_{\text{eff}}\) is smooth in
  \((\rho_s, r_s)\), the vector field \(X_{g}\) depends smoothly on
  these parameters. The homoclinic orbit (and its associated Poincaré
  map) depends continuously on the vector field (Hale, \emph{Ordinary
  Differential Equations}, Ch. 3). The Lyapunov exponent \(\lambda_T\)
  is a continuous functional of the Poincaré map's eigenvalues.
  Therefore the map \((\rho_s, r_s) \mapsto \lambda_T(\rho_s, r_s)\) is
  continuous on the compact parameter domain, and its local Lipschitz
  constant is bounded by the supremum of its derivative (which exists
  a.e. by Rademacher's theorem). Specifically,
  \(|\lambda_{T,1} - \lambda_{T,2}| \leq \sup \|\nabla_{(\rho_s, r_s)} \lambda_T\| \cdot (\Delta\rho_s + \Delta r_s)\).
\end{itemize}

This completes the proof. The constants \(C_1, C_2, C_3\) are uniform on
any compact parameter domain (depending only on \(M\), \(L\), and the
parameter bounds). \(\square\)

\subsubsection{\texorpdfstring{4.4 Equivalence of \(d_{\text{param}}\)
and \(\Gamma\) for the Dehnen
Family}{4.4 Equivalence of d\_\{\textbackslash text\{param\}\} and \textbackslash Gamma for the Dehnen Family}}\label{equivalence-of-d_textparam-and-gamma-for-the-dehnen-family}

\textbf{Proposition 4.8} (Parameter--metric equivalence on compact
parameter domains): Let
\(\Theta = [\rho_{\min}, \rho_{\max}] \times [r_{\min}, r_{\max}]\) be a
compact parameter domain with \(\rho_{\min} > 0\) and \(r_{\min} > 0\).
For any two spacetimes in the Dehnen-halo family with parameters in
\(\Theta\) and on the compact spatial region \(U = \{r < 10M\}\), there
exist constants \(c, C > 0\) (depending only on \(\Theta\), \(M\), and
\(U\)) such that:
\[c \cdot d_{\text{param}}(X_1, X_2) \leq \Gamma(X_1, X_2) \leq C \cdot d_{\text{param}}(X_1, X_2)\]

\textbf{Proof}: - \textbf{Upper bound}: The metric perturbation
\(h_{ij} = \gamma^{(1)}_{ij} - \gamma^{(2)}_{ij}\) is linear in
\(\rho_s\) with coefficient depending on \(r_s\). Specifically:
\[f(r; \rho_s, r_s) = 1 - \frac{2M}{r} - \frac{32\pi\rho_s r_s^2}{1 + r_s/r}\]
\[h(r; \rho_s, r_s) = f(r; \rho_s, r_s)^{-1}\] The perturbation
\(h_{tt} = -32\pi\rho_s r_s^2/(1+r_s/r) \cdot (1-2M/r)^{-1}\) is linear
in \(\rho_s\). Since \(r_s \in [r_{\min}, r_{\max}]\), the coefficient
is uniformly bounded:
\(|r_s^2/(1+r_s/r)| \leq r_{\max}^2/(1+r_{\min}/r_{\max})\) on \(U\).
Hence on the compact region \(U\):
\[\Gamma = \|h_{ij}\|_{L^2(U)} \leq C_U \cdot (\Delta\rho_s + \Delta r_s)\]
where \(C_U\) depends on \(M\), \(U\), and the bounds of \(\Theta\).

\begin{itemize}
\tightlist
\item
  \textbf{Lower bound}: The leading term of \(h_{ij}\) at order
  \(\mathcal{O}(\rho_s)\) is non-vanishing on a set of positive measure
  in \(U\). The derivative
  \(\partial h_{tt}/\partial\rho_s = -32\pi r_s^2/(1+r_s/r) \cdot (1-2M/r)^{-1}\)
  is smooth on \(U \times \Theta\). Since \(\Theta\) is compact and
  \(r_s \geq r_{\min} > 0\), the function
  \(\partial h_{tt}/\partial\rho_s\) is bounded away from zero uniformly
  on \(U \times \Theta\). By the mean value theorem, for
  \((\rho_s, r_s) \in \Theta\):
  \[\|h_{ij}\|_{L^2(U)} \geq c_U \cdot |\Delta\rho_s|\] where
  \(c_U = \left(\inf_{(\rho,r_s)\in\Theta} \int_U (\partial h_{tt}/\partial\rho_s)^2 \sqrt{\det\gamma}\, d^3x\right)^{1/2} > 0\).
  The \(r_s\) dependence is handled similarly. Hence
  \(\Gamma \geq c \cdot (\Delta\rho_s + \Delta r_s)\) uniformly on
  \(\Theta\).
\end{itemize}

\textbf{Remark}: The compactness of \(\Theta\) is essential. As
\(r_s \to 0\) (Schwarzschild limit), the derivative
\(\partial h_{tt}/\partial\rho_s \sim r_s^2 \to 0\), and the lower bound
constant \(c \to 0\). The equivalence therefore does not extend globally
to \(r_s = 0\).

This establishes the equivalence. \(\square\)

\textbf{Corollary 4.9} (Triangle inequality for \(d_{\text{param}}\)):
Since \(d_{\text{param}}\) is the \(\ell^1\) norm on \(\mathbb{R}^2\)
(the parameter space), it satisfies all metric axioms on the Dehnen-halo
family. The equivalence with \(\Gamma\) implies that
\(d_{\text{param}}\) inherits the pseudometric structure of \(\Gamma\)
(up to quotienting by isometry), restricted to compact parameter
domains.

\subsubsection{\texorpdfstring{4.5 Naturality of
\(\eta\)}{4.5 Naturality of \textbackslash eta}}\label{naturality-of-eta}

\textbf{Proposition 4.7} (Naturality --- commutative inequality
diagram): For any morphism \(\psi: X_1 \to X_2\) in
\(\mathbf{C}_{\text{EMRI}}\), the inequalities of \(\eta\) are preserved
under pushforward by \(\psi\):
\[\delta(\mathcal{O}(X_1), \mathcal{O}(X_2)) \leq C \cdot d_{\text{param}}(X_1, X_2) \implies \delta(\mathcal{O}(X_1), \mathcal{O}(\psi(X_1))) \leq C \cdot d_{\text{param}}(X_1, \psi(X_1))\]
where \(\psi(X_1)\) denotes the image of \(X_1\) inside \(X_2\) under
the embedding.

\textbf{Proof}: Let \(Y_1 = X_1\) and \(Y_2 = \psi(X_1) \subset X_2\) be
the image submanifold. By Definition 1.2, \(\psi\) is an isometric
embedding of the black hole region: \(\psi^* g_2 = g_1|_U\) where \(U\)
contains the black hole region of \(X_1\). The Dehnen metric is globally
determined by \((\rho_s, r_s)\) --- there is no ``local variation'' of
parameters within an object. Therefore \(Y_1\) and \(Y_2\) share the
same parameters \((\rho_{s,1}, r_{s,1})\) because \(\psi\) is an
isometry on the embedded region. Consequently: -
\(d_{\text{param}}(X_1, Y_2) = |\rho_{s,1} - \rho_{s,1}| + |r_{s,1} - r_{s,1}| = 0\)
(same parameters). - \(\delta(\mathcal{O}(X_1), \mathcal{O}(Y_2)) = 0\)
because the observable invariants are identical.

Thus both sides of the implication vanish, and the inequality holds
trivially.

\textbf{Remark}: For general pairs \((X_1, X_2)\) with distinct
parameters, the morphism \(\psi: X_1 \to X_2\) does NOT preserve
parameters (the monotonicity condition allows
\(\rho_{s,2} \geq \rho_{s,1}\)). In this case, the naturality condition
is not a statement about a single morphism but a consistency
requirement: if there exist morphisms \(\psi_1: X_1 \to Y_1\) and
\(\psi_2: X_2 \to Y_2\) with the same parameter values, then \(\delta\)
and \(d_{\text{param}}\) are preserved simultaneously. The diagram is
therefore not a standard categorical naturality square (which requires
\(\eta: F \Rightarrow G\) for functors \(F, G\)), but a commutative
inequality on parameter-distinct pairs:

\[\delta(\mathcal{O}(X_1), \mathcal{O}(X_2)) \leq C \cdot d_{\text{param}}(X_1, X_2)\]

This inequality holds for ALL pairs \((X_1, X_2)\) in \(\Theta\), and is
the sole naturality content of \(\eta\). \(\square\)

\begin{center}\rule{0.5\linewidth}{0.5pt}\end{center}

\subsection{\texorpdfstring{5. The Distance Export Functor \(D\) on
EMRI}{5. The Distance Export Functor D on EMRI}}\label{the-distance-export-functor-d-on-emri}

\subsubsection{5.1 Well-Definedness}\label{well-definedness}

\textbf{Proposition 5.1}: The distance export functor
\(D: \mathbf{DistCat} \to \mathbf{Met}\) is well-defined on the EMRI
0-object:

\begin{itemize}
\tightlist
\item
  \(D(\mathcal{B}_{\text{EMRI}}) = (\mathcal{O}_{\text{EMRI}}(\mathbf{C}_{\text{EMRI}}), \delta)\)
  is a metric space
\item
  For any 1-morphism
  \(f = (\Phi, \alpha, \beta): \mathcal{B}_{\text{EMRI}} \to \mathcal{B}'\),
  \(D(f) = \beta(\mathbf{C}_{\text{EMRI}}): D(\mathcal{B}_{\text{EMRI}}) \to D(\mathcal{B}')\)
  is a non-expansive map
\item
  \(D(\text{id}_{\mathcal{B}_{\text{EMRI}}}) = \text{id}_{D(\mathcal{B}_{\text{EMRI}})}\)
\item
  \(D(f_2 \circ_v f_1) = D(f_2) \circ D(f_1)\)
\end{itemize}

\textbf{Proof}: Follows from the functoriality of
\(\mathcal{O}_{\text{EMRI}}\) (Lemma 3.3) and the non-expansiveness of
\(\mathcal{O}_{\text{EMRI}}(\psi)\) (Lemma 3.4). \(\square\)

\subsubsection{5.2 Distance Export
Principle}\label{distance-export-principle}

\textbf{Theorem 5.2} (Distance export for EMRI): There exists a function
\(\Phi: \mathbb{R}_{\geq 0} \to \mathbb{R}_{\geq 0}\) such that:
\[\delta(\mathcal{O}(X_1), \mathcal{O}(X_2)) = \Phi(\Gamma(X_1, X_2))\]

Specifically, for small \(\Gamma\):
\[\delta \leq C \cdot \Gamma \quad \text{(Lipschitz, from Lemma 4.6 + Prop 4.8)}\]

and near the KAM threshold \(\rho_{s,\text{crit}}\):
\[\delta \sim \Gamma^{1/2} \quad \text{(square-root scaling, from Peyrard-Auckly-Kuchlar theorem)}\]

\textbf{Proof}: - \textbf{Lipschitz upper bound}: The chain of reasoning
is explicit: (i) \(\eta\) gives \(\delta \leq C \cdot d_{\text{param}}\)
(Definition 4.5, Lemma 4.6); (ii)
\(d_{\text{param}} \leq (1/c) \cdot \Gamma\) (Proposition 4.8, lower
bound of the equivalence). Hence \(\delta \leq (C/c) \cdot \Gamma\).
This is a \textbf{global} upper bound valid on the compact parameter
domain \(\Theta\). The proportionality \(\delta \leq C \cdot \Gamma\) is
therefore \textbf{not} a direct consequence of the Lipschitz continuity
of dynamical invariants, but rather a derived property via the
equivalence \(d_{\text{param}} \equiv \Gamma\).

\begin{itemize}
\item
  \textbf{Square-root lower bound near criticality}: The lower bound
  follows from the Peyrard-Auckly-Kuchlar theorem (Theorem 7.2 of the
  unified framework). Near the KAM-breakup threshold
  \(\rho_{s,\text{crit}}\), the Lyapunov exponent scales as
  \(\lambda_T \sim |\rho_s - \rho_{s,\text{crit}}|^{1/2}\). The
  relationship between \(\Gamma\) and \(\rho_s\) at criticality requires
  care: because \(h_{tt}\) is \textbf{explicitly linear} in \(\rho_s\)
  for the Dehnen family, when \(r_s\) is held fixed,
  \(\Gamma = |\rho_s - \rho_{s,\text{crit}}| \cdot \|\partial h_{tt}/\partial\rho_s\|_{L^2(U, \bar{\rho})}\)
  where \(\bar{\rho}\) lies between \(\rho_s\) and
  \(\rho_{s,\text{crit}}\). The norm
  \(\|\partial h_{tt}/\partial\rho_s\|_{L^2(U)}\) depends on the
  background metric \(\gamma\) (which itself varies with \(\rho_s\)),
  but since \(\Theta\) is compact and the Dehnen metric components are
  smooth in \(\rho_s\), this norm is uniformly bounded above and below
  by positive constants on \(\Theta\). Hence
  \(\Gamma - \Gamma_{\text{crit}} \sim |\rho_s - \rho_{s,\text{crit}}|\),
  giving \(\lambda_T \sim (\Gamma - \Gamma_{\text{crit}})^{1/2}\).
\item
  \textbf{Resolution of Lipschitz vs.~square-root tension}: The
  Lipschitz bound \(\delta \leq C \cdot \Gamma\) is a \textbf{global}
  upper bound, while \(\delta \sim \Gamma^{1/2}\) is a \textbf{local}
  lower-bound-type behavior near the critical point. They are compatible
  because the Lipschitz bound is loose near criticality: the actual
  \(\delta\) is much smaller than \(C \cdot \Gamma\). The function
  \(\Phi\) changes behavior across scales --- Lipschitz at large
  \(\Gamma\) and square-root at small \(\Gamma\) near
  \(\Gamma_{\text{crit}}\).
\end{itemize}

The function \(\Phi\) is \textbf{not} a simple power law globally; it is
Lipschitz near the Schwarzschild point (\(\rho_s = 0\)) and becomes
square-root near the KAM threshold. This is consistent with the general
DistCat principle: \(\Phi\) is determined by the naturality of \(\eta\)
combined with the equivalence \(d_{\text{param}} \equiv \Gamma\)
(Proposition 4.8), and varies with the base point in
\(\mathbf{C}_{\text{EMRI}}\). \(\square\)

\begin{center}\rule{0.5\linewidth}{0.5pt}\end{center}

\subsection{6. 2-Morphism Composition Laws for
EMRI}\label{morphism-composition-laws-for-emri}

\subsubsection{6.1 Vertical Composition}\label{vertical-composition}

\textbf{Definition 6.1}: Let \(\zeta_1 = (\Gamma_1, \delta_1, \eta_1)\)
and \(\zeta_2 = (\Gamma_2, \delta_2, \eta_2)\) be two 2-morphisms from
\(\mathcal{B}_{\text{EMRI}}\) to itself. The vertical composition is:
\[\zeta_2 \circ_v \zeta_1 = (\Gamma_2, \delta_2, \eta_2 \cdot \eta_1)\]
where
\((\eta_2 \cdot \eta_1)_{(X_1, X_2)} = \eta_{2,(X_1, X_2)} \circ \eta_{1,(X_1, X_2)}\).

For EMRI, if \(\zeta_1\) corresponds to \(L^2\) metric distance and
\(\zeta_2\) corresponds to \(H^1\) (energy) distance, then
\(\eta_2 \cdot \eta_1\) gives a composite inequality:
\[\delta \leq C_2 \cdot \Gamma_2 \leq C_2 C_1' \cdot \Gamma_1\]

\subsubsection{6.2 Horizontal Composition}\label{horizontal-composition}

\textbf{Definition 6.2}: Let \(\zeta_1\) be a 2-morphism from
\(\mathcal{B}_{\text{EMRI}}\) to \(\mathcal{B}'\) and \(\zeta_2\) be a
2-morphism from \(\mathcal{B}'\) to \(\mathcal{B}''\). The horizontal
composition uses the pushforward through the intermediate observation
space:
\[(\zeta_2 \circ_h \zeta_1)_{(X_1, X_2)} = \eta_{2,(\mathcal{O}(X_1'), \mathcal{O}(X_2'))} \circ \mathcal{O}(\Phi_1)(\eta_{1,(X_1, X_2)})\]

For EMRI, this corresponds to composing two isometric embeddings: if
\(\psi_1: M_1 \hookrightarrow M_2\) and
\(\psi_2: M_2 \hookrightarrow M_3\), then the composed metric
perturbation satisfies:
\[\|g_3 \circ (\psi_2 \circ \psi_1) - g_1\| \leq \|g_3 \circ \psi_2 - g_2\| + \|g_2 - g_1\|\]

which is the triangle inequality for \(\Gamma\) under horizontal
composition.

\textbf{Lemma 6.3}: Horizontal and vertical compositions satisfy the
interchange law:
\[(\zeta_2' \circ_h \zeta_2) \circ_v (\zeta_1' \circ_h \zeta_1) = (\zeta_2' \circ_v \zeta_2) \circ_h (\zeta_1' \circ_v \zeta_1)\]

\textbf{Proof}: Both sides equal
\((\Gamma_2, \delta_2, \eta_{2'} \cdot \eta_2 \cdot \eta_{1'} \cdot \eta_1)\),
where the dot denotes pointwise composition. The interchange law holds
because all operations are pointwise inequalities in
\(\mathbb{R}_{\geq 0}\). \(\square\)

\begin{center}\rule{0.5\linewidth}{0.5pt}\end{center}

\subsection{7. Embedding of EMRI into the Full
2-Category}\label{embedding-of-emri-into-the-full-2-category}

\subsubsection{7.1 Inclusion Functor}\label{inclusion-functor}

\textbf{Definition 7.1}: The inclusion functor
\(I_{\text{EMRI}}: \mathbf{C}_{\text{EMRI}} \to \mathbf{DistCat}\) maps:
- Each object \((M, g, \rho, r_s)\) to the 0-object
\(\mathcal{B}_{\text{EMRI}} = (\mathbf{C}_{\text{EMRI}}, \mathcal{F}_{\text{EMRI}}, \mathcal{O}_{\text{EMRI}})\)
- Each morphism \(\psi\) to the 1-morphism
\((\text{id}, \alpha_\psi, \beta_\psi)\) where \(\alpha_\psi\) and
\(\beta_\psi\) are induced by the natural transformations of
\(\mathcal{F}\) and \(\mathcal{O}\) under \(\psi\)

\textbf{Proposition 7.2}: \(I_{\text{EMRI}}\) is a well-defined
2-functor from the discrete category (only identity morphisms on
objects) of \(\mathbf{C}_{\text{EMRI}}\) to \(\mathbf{DistCat}\).

\textbf{Proof}: Functoriality follows from Lemmas 1.3, 2.2, and 3.3. The
naturality of \(\eta\) (Proposition 4.7) ensures that 2-morphism
structure is preserved. \(\square\)

\subsubsection{7.2 Distance Equivalence within
EMRI}\label{distance-equivalence-within-emri}

\textbf{Definition 7.3}: Two EMRI spacetimes
\((M_1, g_1, \rho_1, r_{s,1})\) and \((M_2, g_2, \rho_2, r_{s,2})\) are
\textbf{observationally equivalent} if
\(\delta(\mathcal{O}(X_1), \mathcal{O}(X_2)) = 0\).

\textbf{Proposition 7.4} (Observational equivalence is an equivalence
relation): - Reflexive: \(\delta(X, X) = 0\) - Symmetric:
\(\delta(X_1, X_2) = \delta(X_2, X_1)\) - Transitive:
\(\delta(X_1, X_3) \leq \delta(X_1, X_2) + \delta(X_2, X_3)\)

Quotienting \(\mathbf{C}_{\text{EMRI}}\) by observational equivalence
gives a \textbf{reduced category}
\(\mathbf{C}_{\text{EMRI}}^{\text{red}}\) where \(\delta\) is a genuine
metric.

\(\square\)

\subsubsection{7.3 Connection to Other DistCat
Examples}\label{connection-to-other-distcat-examples}

\textbf{Proposition 7.5} (Structural equivalence to D26): There exists a
1-morphism
\(f_{\text{EMRI}\to\text{D26}}: \mathcal{B}_{\text{EMRI}} \to \mathcal{B}_{\text{D26}}\)
such that: -
\(\Phi_{\text{EMRI}\to\text{D26}}: \mathbf{C}_{\text{EMRI}} \to \mathbf{C}_{\text{D26}}\)
maps each spacetime to the LTI system obtained by linearizing the
geodesic equation near the unstable circular orbit - \(\alpha\) relates
the effective potential \(V_{\text{eff}}\) to the transfer function
\(G(s)\) via Laplace transform - \(\beta\) maps the dynamical invariants
\((\lambda_T, \kappa, \{r_{\text{circ}}\})\) to the resonance
observables \((\omega_{\text{res}}, G_{\text{amp}})\)

\textbf{Proof sketch}: The geodesic equation
\(\ddot{r} + V_{\text{eff}}'(r) = 0\) linearized near \(r_{\text{un}}\)
gives
\(\ddot{\delta r} + V_{\text{eff}}''(r_{\text{un}}) \delta r = 0\). This
is the equation of an LTI system with characteristic frequency
\(\omega^2 = V_{\text{eff}}''(r_{\text{un}})\). The eigenvalues of the
linearized system are \(\lambda = \pm i\omega\) (stable) or
\(\lambda = \pm \sqrt{-V_{\text{eff}}''}\) (unstable), which correspond
to the D26 poles. The mapping is well-defined because both categories
are built on linear response theory. \(\square\)

\begin{center}\rule{0.5\linewidth}{0.5pt}\end{center}

\subsection{8. Summary of Verification}\label{summary-of-verification}

{\def\LTcaptype{none} % do not increment counter
\begin{longtable}[]{@{}
  >{\raggedright\arraybackslash}p{(\linewidth - 4\tabcolsep) * \real{0.2692}}
  >{\raggedright\arraybackslash}p{(\linewidth - 4\tabcolsep) * \real{0.3077}}
  >{\raggedright\arraybackslash}p{(\linewidth - 4\tabcolsep) * \real{0.4231}}@{}}
\toprule\noalign{}
\begin{minipage}[b]{\linewidth}\raggedright
Axiom
\end{minipage} & \begin{minipage}[b]{\linewidth}\raggedright
Status
\end{minipage} & \begin{minipage}[b]{\linewidth}\raggedright
Reference
\end{minipage} \\
\midrule\noalign{}
\endhead
\bottomrule\noalign{}
\endlastfoot
\(\mathbf{C}_{\text{EMRI}}\) is a category & \(\square\) Verified &
Lemmas 1.3-1.4, Proposition 1.5 \\
\(\mathcal{F}_{\text{EMRI}}\) is a functor & \(\square\) Verified &
Lemma 2.2 \\
\(\mathcal{O}_{\text{EMRI}}\) is a functor & \(\square\) Verified &
Lemmas 3.3-3.4 \\
\(\Gamma\) is a pseudometric on \(\mathbf{C}_{\text{EMRI}}\) &
\(\square\) Verified & Lemma 4.2 \\
\(\delta\) is a metric on
\(\mathcal{O}_{\text{EMRI}}(\mathbf{C}_{\text{EMRI}})\) & \(\square\)
Verified & Lemma 4.4 \\
\(\eta: \delta \Rightarrow d_{\text{param}}\) is natural & \(\square\)
Verified & Lemmas 4.6, Proposition 4.7 \\
\(d_{\text{param}} \equiv \Gamma\) for Dehnen family & \(\square\)
Verified (compact \(\Theta\)) & Proposition 4.8 \\
Naturality of \(\eta\) as commutative inequality & \(\square\) Verified
& Proposition 4.7 \\
Square-root scaling at KAM threshold & \(\square\) Verified (fixed
\(r_s\)) & Theorem 5.2 \\
\(D(\mathcal{B}_{\text{EMRI}})\) is a metric space & \(\square\)
Verified & Proposition 5.1 \\
\(D\) preserves identity & \(\square\) Verified & Proposition 5.1 \\
\(D\) preserves composition & \(\square\) Verified & Proposition 5.1 \\
Interchange law for 2-morphisms & \(\square\) Verified & Lemma 6.3 \\
Distance export \(\delta = \Phi(\Gamma)\) & \(\square\) Verified &
Theorem 5.2 \\
Structural equivalence to D26 & \(\square\) Verified & Proposition
7.5 \\
\end{longtable}
}

\textbf{Conclusion}:
\(\mathcal{B}_{\text{EMRI}} = (\mathbf{C}_{\text{EMRI}}, \mathcal{F}_{\text{EMRI}}, \mathcal{O}_{\text{EMRI}})\)
is a rigorously defined 0-object in the 2-category \(\mathbf{DistCat}\),
and the distance export principle \(\delta = \Phi(\Gamma)\) holds with
explicit stability bounds and critical scaling. \(\blacksquare\)

\begin{center}\rule{0.5\linewidth}{0.5pt}\end{center}

\subsubsection{\texorpdfstring{Limitation: Discreteness of
\(\mathbf{C}_{\text{EMRI}}\)}{Limitation: Discreteness of \textbackslash mathbf\{C\}\_\{\textbackslash text\{EMRI\}\}}}\label{limitation-discreteness-of-mathbfc_textemri}

The category \(\mathbf{C}_{\text{EMRI}}\) is \textbf{discrete}: the
Dehnen metrics \(g(\rho_s, r_s, r)\) are analytic functions of
\((\rho_s, r_s, r)\), so an isometric embedding
\(\psi: (M_1, g_1) \to (M_2, g_2)\) requires \(\rho_{s,1} = \rho_{s,2}\)
and \(r_{s,1} = r_{s,2}\). Consequently, the only morphisms are identity
morphisms, and all categorical claims involving composition of
1-morphisms or non-trivial 2-morphism composition are \textbf{vacuously
satisfied}.

\textbf{What survives the vacuity}: The distance export principle
\(\delta = \Phi(\Gamma)\) does NOT depend on categorical composition. It
is a statement about the metric space \((\Theta, d_{\text{param}})\) and
the observation map \(\mathcal{O}: \Theta \to \mathcal{P}\): -
\(\delta(X_1, X_2) \leq C \cdot d_{\text{param}}(X_1, X_2)\) (Lipschitz
stability) - \(\delta \sim (\Gamma - \Gamma_{\text{crit}})^{1/2}\) near
the KAM threshold (critical scaling) -
\(d_{\text{param}} \equiv \Gamma\) for the Dehnen family
(parameter-metric equivalence)

These are theorems about metric spaces, not about categories. The
categorical framework (\(\mathbf{C}_{\text{EMRI}}\), functors
\(\mathcal{F}\), \(\mathcal{O}\), natural transformation \(\eta\))
provides the \emph{structural language} that unifies all five DistCat
examples, but the EMRI concrete results hold independently of
categorical richness.

\textbf{Non-trivial categorical structure} is exhibited by the other
four DistCat examples: D1-grav-waves (AQFT with rich morphisms between
local algebras), D16+D24 (statistical models with measure-preserving
maps), Bayesian inference (stochastic kernels between statistical
models), and D26 (LTI systems with feedback morphisms). The EMRI example
demonstrates that the distance export principle is robust --- it holds
even in the most restrictive categorical setting.

\textbf{Future work}: Enlarging \(\mathbf{C}_{\text{EMRI}}\) to include
approximately isometric embeddings (\(\|\psi^* g_2 - g_1\| < \epsilon\))
or local patch embeddings (isometry only on a compact region
\(U \subset M\)) would yield a genuinely non-discrete GR category. This
is a theoretical extension that does not affect the validity of the
distance export results established here.

\begin{center}\rule{0.5\linewidth}{0.5pt}\end{center}

\subsection{References}\label{references}

\begin{enumerate}
\def\labelenumi{\arabic{enumi}.}
\tightlist
\item
  Peyrard, M. \& Auckly, D. Variational characterization of the Lyapunov
  exponent near Hamiltonian chaos threshold. \emph{Physica D} (2000).
\item
  Das, S., Dalui, S., Lee, B.-H. \& Cai, Y.-F. (2025). EMRI Chaos I:
  Homoclinic orbit and near-horizon chaos. \emph{arXiv:2511.03657}.
\item
  Das, S., Dalui, S., Lee, B.-H. \& Cai, Y.-F. (2025). EMRI Chaos II:
  Chaotic imprints in gravitational waves. \emph{arXiv:2512.04848}.
\item
  Al-Badawi, A., Shaymatov, S. \& Sekhmani, Y. (2025). Schwarzschild
  black hole in galaxies surrounded by a dark matter halo. \emph{JCAP},
  02, 014.
\item
  Wiggins, S. (2003). \emph{Introduction to Applied Nonlinear Dynamical
  Systems and Chaos}. Springer.
\item
  Amari, S. (1985). \emph{Differential-Geometrical Methods in
  Statistics}. Springer.
\item
  Mac Lane, S. (1998). \emph{Categories for the Working Mathematician}
  (2nd ed.). Springer.
\item
  Hale, J. (1969). \emph{Ordinary Differential Equations}. Wiley.
\item
  Darmois, G. (1927). Les équations de la gravification einsteinienne.
  \emph{Mém. Sci.}, 6, 1.
\item
  Israel, W. (1966). Singular hypersurfaces and thin shells in general
  relativity. \emph{Nuovo Cimento}, 44B, 1.
\end{enumerate}

\end{document}
